
The fundamental goal of LBMPC is to improve performance
by learning unknown system dynamics $g(x, u)$ online
while maintaining the stability and safety guarantees of robust MPC.

\subsection{Regression Techniques for Model Learning}

Before controlling the system, one must quantify the uncertainty.

\textbf{Parametric vs. Non-parametric}
Parametric models (like BLR) use a finite set of weights $\theta$.
Non-parametric models (like GPs) use the data points themselves.

\textbf{Robust vs. Stochastic}
Robust methods (Set-Membership) provide hard bounds.
Stochastic methods (BLR, GP) provide probability distributions.

\begin{enumerate}
  \item
    \textbf{Bayesian Linear Regression (BLR)}
    Assumes $y = \Phi(x)^\top \theta + w$.
    It updates a Gaussian posterior
    $p(\theta | \mathcal{D}) = \mathcal{N}(\mu_\theta, \Sigma_\theta)$
    using Bayes' rule.
  \item
    \textbf{Gaussian Processes (GP)}
    A non-parametric method
    defined by a mean $m(x)$ and a kernel $k(x, x')$.
    It provides a mean prediction $\hat{f}(x_*)$
    and a variance $\text{var}[f(x_*)]$.
  \item
    \textbf{Set-Membership (SM) Estimation}
    A robust approach.
    Given a bound on noise $w \in \mathcal{W}$,
    it computes a \textbf{non-falsified set} $\Delta_i$
    of parameters consistent with the data.
    The parameter set $\Omega_i$ is updated via intersection:
    $\Omega_i = \Omega_{i-1} \cap \Delta_i$.
\end{enumerate}

\subsection{Parameter Estimation Methods}

\subsubsection{System Modeling}

For parameter $\theta$ and  with $w_k, v_k$ disturbances/noise:
\[
  x_{k+1} = f(x_k, u_k, \theta) + w_k, \quad y_k = h(x_k, \theta) + v_k
\]
Estimation refines $\theta$ using observed data. Divided into offline
(pre-computed) and online (real-time) approaches.

\textbf{Offline Estimation}
Uses historical data for initial $\theta$ fitting.

Least Squares (LS): Minimizes residual sum:
$$
\hat{\theta} = \arg\min_{\theta} \sum_{i=1}^N \| y_i - \phi_i^T \theta \|^2
$$
where $\phi_i$ is the regressor vector.

Extensions:
Recursive LS (RLS) for incremental updates,
Bayesian methods for uncertainty quantification.

\textbf{Online Estimation}
Adapts $\theta$ during operation, enabling learning.

Kalman Filter (KF)-based:
Treats $\theta$ as augmented state, updating via prediction-correction.
\[
  \hat{\theta}_{k|k} = \hat{\theta}_{k|k-1} + K_k (y_k - \hat{y}_{k|k-1})
\]
where $K_k$ is Kalman gain.

Machine Learning Integration:
Gaussian Processes (GP) or Neural Networks (NN)
for nonparametric estimation,
capturing complex $\theta$ dynamics.
GP models uncertainty as:
\[
  f(\cdot) \sim \mathcal{GP}(m(\cdot), k(\cdot, \cdot))
\]
with mean $m$ and kernel $k$.

\subsection{Robust Performance LBMPC}

Core objective is to \textbf{decouple} performance from safety.

The controller uses two models simultaneously:

\textbf{Nominal Linear Model}
\[
  z_{i+1} = Az_i + Bu_i
\]
Used to ensure safety.
Constraints are tightened relative to this model using
Robust MPC techniques (e.g., tube-based MPC).

\textbf{Learned Nonlinear Model}
\[
  x_{i+1} = Ax_i + Bu_i + \mathcal{O}_k(x_i, ui)
\]
The \textit{Oracle} $\mathcal{O}_k$ represents the learned dynamics.
The MPC \textbf{cost function} is evaluated using
this model to achieve high performance.

\textbf{Result}
Even if the learned model is wrong,
the nominal model ensures the system remains within the \textit{safe} tube.

\subsection{Robust Adaptive MPC (Set-Membership Update)}

While the Performance model can be any black box,
\textbf{Uncertainty Set} $\mathcal{W}$ must be updated carefully
to avoid losing feasibility.

\textbf{Mechanism} As new data arrives, the set $\Omega$ shrinks.

\textbf{Fixed Shape Parameter Set}
To prevent the number of constraints from growing to infinity,
the set $\Omega$ is over-approximated by a polytope of fixed shape,
where only the offsets $h$ are updated via a Linear Program (LP).

\textbf{Guarantee}
Because $\Omega_{i+1} \subseteq \Omega_i$,
the feasible region of the MPC is non-decreasing
($X_{feas}^{i+1} \supseteq X_{feas}^i$),
ensuring \textbf{recursive feasibility}.

\subsection{Stochastic Learning-based MPC (Cautious MPC)}

When using GPs, the uncertainty is propagated through the horizon.

\textbf{Uncertainty Propagation}
Since propagating a Gaussian through a nonlinear GP is hard,
approximations are used:

\textbf{Mean Equivalent}
Treat GP input as deterministic at its mean.

\textbf{Taylor Expansion}
Linearize the GP mean around the current state mean.

\textbf{Exact Moment Matching}
Analytically compute mean and variance (possible for RBF kernels).

\textbf{Cautious MPC}
The controller is \textit{cautious} because
the variance $\Sigma(x)$ appears in the constraints
(chance constraints) and the cost function.

